\documentclass[11pt]{article}
\usepackage[T1]{fontenc}
\usepackage{amsmath,amssymb,amsthm}
\usepackage[margin=1in]{geometry}
\usepackage[hidelinks]{hyperref}
\usepackage{booktabs}

\newtheorem{theorem}{Theorem}
\newtheorem{lemma}[theorem]{Lemma}
\newcommand{\RP}{\mathrm{RP}}
\newcommand{\RN}{\mathrm{RN}}

\title{The Exact Resistance-Positivity Threshold for a Subdivided-Clique Join}
\author{Xiaojian Zeng\\
PBC School of Finance, Tsinghua University, 100190\\
\texttt{xiaojianzeng@foxmail.com}}
\date{}

\begin{document}
\maketitle

\begin{abstract}
For integers $m\geq2$ and $\ell\geq1$, let $H_m$ be obtained from the
complete graph on $\{a,b_1,\ldots,b_m\}$ by subdividing every edge $ab_i$
once, and put $J_{m,\ell}=K_\ell\vee H_m$.  Wang proved that
$J_{m,\ell}$ is not resistance-nonnegative when $m\geq2\ell+4$.
We prove the complementary positive range: $J_{m,\ell}$ is
resistance-positive whenever $m\leq2\ell+3$.  Below the boundary we give
Hamiltonian cycles.  At $m=2\ell+3$ we give seven rational edge-orbit
marginals in the relative interior of the spanning-tree polytope, with every
vertex incident sum strictly below two.  Thus
$J_{m,\ell}$ is resistance-positive exactly when $m\leq2\ell+3$.
\end{abstract}

\section{Context and statement}

All graphs are finite, simple, undirected, and connected.  Given positive
edge conductances $c_e$, write $r_e(c)$ for the relative resistance of an
edge $e$.  The resistance curvature at a vertex is
\[
 p_v(c)=1-\frac12\sum_{e\ni v}r_e(c).
\]
A graph is resistance-positive (\(\RP\)) if some positive conductances make
all $p_v(c)$ positive, and resistance-nonnegative (\(\RN\)) if they make all
of them nonnegative.  Devriendt introduced these graph classes and proved a
spanning-tree-distribution characterization of them
\cite{Devriendt2026}.  Recent work has developed further structural and
polyhedral aspects of \(\RP\) and \(\RN\) graphs
\cite{AgrahariEtAl2026}.

Let
\[
 B=\{b_1,\ldots,b_m\},\qquad X=\{x_1,\ldots,x_m\},
 \qquad U=\{u_1,\ldots,u_\ell\}.
\]
The graph $H_m$ has vertex set $\{a\}\cup B\cup X$ and edge set
\[
 \{b_ib_j:1\leq i<j\leq m\}\cup
 \{ax_i,x_ib_i:1\leq i\leq m\}.
\]
The graph $J_{m,\ell}=K_\ell\vee H_m$ additionally has every edge within
$U$ and every edge between $U$ and $V(H_m)$.  Wang introduced this join
family, proved $\tau(J_{m,\ell})=1+\ell/m$, and proved that it is not \(\RN\)
for $m\geq2\ell+4$ \cite[Propositions 15--16]{Wang2026}.
Our contribution is the complementary assertion in the following theorem.

\begin{theorem}\label{thm:main}
For all integers $m\geq2$ and $\ell\geq1$,
\[
 J_{m,\ell}\text{ is \(\RP\)}\quad\Longleftrightarrow\quad
 m\leq2\ell+3.
\]
\end{theorem}

The negative implication is Wang's Proposition~16.  We prove the positive
implication.  The boundary $m=2\ell+3$ is the only part not covered by the
Hamiltonian constructions below.

\section{The polyhedral criterion}

For a connected graph $G=(V,E)$, let $P(G)$ be the convex hull in
$\mathbb R^E$ of the indicator vectors of its spanning trees.  For
$x\in\mathbb R^E$ and $F\subseteq E$, write $x(F)=\sum_{e\in F}x_e$; for
$S\subseteq V$, let $E[S]$ be the edges with both endpoints in $S$ and let
$E(v)$ be the edges incident with $v$.  The graphic-matroid base-polytope
description goes back to Edmonds \cite{Edmonds1971}.  Combined with
Devriendt's spanning-tree characterization, it gives the following criterion;
see also the explicit resistance-curvature formulation in
\cite{Devriendt2026,AgrahariEtAl2026}.

\begin{lemma}\label{lem:criterion}
Let $G=(V,E)$ be $2$-connected.  If $x\in\mathbb R^E$ satisfies
\begin{align*}
 x_e&>0 &&(e\in E),\\
 x(E)&=|V|-1,\\
 x(E[S])&<|S|-1 &&(2\leq |S|\leq |V|-1),\\
 x(E(v))&<2 &&(v\in V),
\end{align*}
then $G$ is \(\RP\).
\end{lemma}

\begin{proof}
The first three lines say precisely that $x$ lies in the relative interior
of $P(G)$.  Hence $x$ is the edge-marginal vector of a positive distribution
on all spanning trees.  Under that distribution, the expected degree of
$v$ is $x(E(v))<2$.  Devriendt's characterization now gives \(\RP\).
\end{proof}

\section{Hamiltonian cycles below the boundary}

Put $L=\ell$ and consider the cyclic list
\[
 c_0=a,c_1=u_1,\ldots,c_L=u_L,c_{L+1}=c_0.
\]
It has $L+1$ cyclic positions.  When $L=1$, the two positions are the two
oriented arcs between $a$ and $u_1$.

Suppose first that $m=2r$.  Since $m\leq2L+2$, we have $r\leq L+1$.
For $j=1,\ldots,r$, replace the position $c_{j-1}c_j$ by
\begin{equation}\label{eq:pairblock}
 c_{j-1},x_{2j-1},b_{2j-1},b_{2j},x_{2j},c_j,
\end{equation}
and retain all other positions.  Every consecutive pair is adjacent and
every vertex occurs exactly once, so the resulting cyclic list is a
Hamiltonian cycle.

Now suppose that $m=2r+1$.  Then $r\leq L$.  If $L=1$, necessarily $m=3$,
and a Hamiltonian cycle is
\[
 a,x_1,b_1,b_2,b_3,x_3,u_1,x_2,a.
\]
If $L\geq2$, replace $c_0c_1$ by
\[
 c_0,x_1,b_1,b_2,b_3,x_3,c_1,
\]
replace $c_1c_2$ by $c_1,x_2,c_2$.  Put $r-1$ further blocks of the form
\eqref{eq:pairblock} into any $r-1$ of the remaining $L-1$ positions,
using the index pairs
\[
 (4,5),(6,7),\ldots,(2r,2r+1).
\]
Again this is a Hamiltonian cycle.

For completeness, every Hamiltonian graph is \(\RP\).  If its Hamiltonian
cycle has $n$ vertices, uniformly delete one cycle edge.  The resulting
random spanning path has expected degree $2-2/n$ at each vertex.  Mixing
this distribution with a sufficiently small positive multiple of the
uniform distribution on all spanning trees gives full support while
preserving every strict expected-degree bound.  Thus the preceding cycles
prove \(\RP\) for $m\leq2L+2$.

\section{An exact boundary certificate}

Assume $m=2L+3$ and set
\[
 A=3L+4,\qquad D=(L+2)(3L+4).
\]
Assign a common marginal to every edge in each of the seven edge orbits as
follows:
\begin{equation}\label{eq:orbits}
\begin{array}{c|ccccccc}
\text{orbit}&BB&aX&b_ix_i&UU&Ua&UB&UX\\ \hline
x_e&\dfrac{3L^2+9L+8}{2(L+1)D}&
\dfrac{3L+5}{D}&1-\dfrac1D&
\dfrac1D&\dfrac1D&\dfrac1D&\dfrac1{L+2}.
\end{array}
\end{equation}
Here $BB$ denotes an edge between two vertices of $B$, and the other orbit
names have the analogous meaning.  All seven numbers are positive.  Counting
the edges in each orbit gives
\begin{equation}\label{eq:total}
 x(E)=2m+L=|V(J_{m,L})|-1.
\end{equation}
The incident sums at $a$, at any $b_i$, at any $x_i$, and at any $u_j$ are,
respectively,
\begin{equation}\label{eq:incidence}
 2-\frac1D,\qquad 2-\frac1D,\qquad
 2-\frac1{L+2},\qquad 2-\frac1D.
\end{equation}
It remains to verify all rank inequalities strictly.

Let $S$ be a proper vertex subset with $|S|\geq2$, and encode it by
\[
 \alpha=\boldsymbol 1_{a\in S},\quad k=|S\cap B|,\quad
 h=|S\cap X|,\quad z=|S\cap U|,
\]
together with the number $d$ of indices $i$ for which both $b_i,x_i\in S$.
Then
\begin{align}
x(E[S])={}&\frac{3L^2+9L+8}{2(L+1)D}\binom{k}{2}
 +\frac{3L+5}{D}\alpha h+\left(1-\frac1D\right)d \notag\\
&+\frac1D\left[\binom{z}{2}+z\alpha+zk\right]
 +\frac1{L+2}zh.                                      \label{eq:internal}
\end{align}
For fixed $\alpha,k,h,z$, the right-hand side is largest at
$d=\min(k,h)$, and that value is attainable.  Put
\[
 P=\frac{3L^2+9L+8}{2(L+1)},\qquad
 \Phi=D\bigl(|S|-1-x(E[S])\bigr)
\]
after this worst-case substitution.  If $k\leq h$, then
\begin{align}
\Phi={}&D(\alpha+h+z-1)-P\binom{k}{2}-(A+1)\alpha h
-\binom{z}{2}-z\alpha-zk-Azh+k,                    \label{eq:phi1}
\end{align}
whereas if $h\leq k$, then
\begin{align}
\Phi={}&D(\alpha+k+z-1)-P\binom{k}{2}-(A+1)\alpha h
-\binom{z}{2}-z\alpha-zk-Azh+h.                    \label{eq:phi2}
\end{align}

We give the endpoint reduction, including the size and full-set exceptions.
In \eqref{eq:phi1}, the coefficient of $h$ is
\[
 D-(A+1)\alpha-Az\geq A-1>0,
\]
so the minimum occurs at $h=k$, except where the lower constraint $|S|=2$
forces an adjacent value.  In \eqref{eq:phi2}, the coefficient of $h$ is
$1-(A+1)\alpha-Az$.  Thus $h=0$ when $\alpha=z=0$, and $h=k$ otherwise,
apart from the point adjacent to the excluded full vertex set.  On the
diagonal $h=k$, and also on the branch $\alpha=z=h=0$, the resulting
expression is concave in $k$.  After taking the interval endpoints in $k$,
it is concave in $z$.  Consequently only endpoint values and the stated
adjacent lattice points remain.  After duplicates are removed, they are
listed below; the row $(0,0,m,0)$ is included although it is dominated by
the displayed $k=2$ endpoint.

\begin{center}
\small
\begin{tabular}{@{}lll@{}}
\toprule
$(\alpha,z,k,h)$ or family & condition & $\Phi$ \\
\midrule
$(0,0,1,1)$ && $1$\\
$(0,0,0,2)$ && $D$\\
$(0,0,2,0)$ && $\dfrac{6L^3+23L^2+27L+8}{2(L+1)}$\\
$(0,0,m,0)$ && $\dfrac{6L^3+25L^2+29L+8}{2}$\\
$(1,0,0,1)$ && $3L^2+7L+3$\\
$(0,1,0,1),(0,1,1,1),(1,0,1,1)$ && $(L+1)(3L+4)$\\
$k=0,z=2,\alpha=0$ & $L\geq2$ & $D-1$\\
$k=0,z=L,\alpha=0$ & $L\geq2$ & $\dfrac{(L-1)(6L^2+19L+16)}2$\\
$k=0,z=1,\alpha=1$ && $D-1$\\
$k=0,z=L,\alpha=1$ && $\dfrac{L(2L+3)(3L+5)}2$\\
$k=h=m,z=0,\alpha=0$ && $\dfrac{(L+1)(L+2)(6L+7)}2$\\
$k=h=m,z=L,\alpha=0$ && $(L+1)(3L+7)$\\
$k=h=m,z=0,\alpha=1$ && $\dfrac{L(2L+3)(3L+5)}2$\\
$k=h=m,z=L-1,\alpha=1$ && $(L+1)(3L+7)$\\
$\alpha=1,z=L;(k,h)=(m-1,m-1),(m,m-1)$
&& $(L+1)(3L+4)$\\
\bottomrule
\end{tabular}
\end{center}

Every entry is at least one for $L\geq1$ (the two indicated rows occur only
for $L\geq2$).  Therefore every proper $S$ with $|S|\geq2$ satisfies
\begin{equation}\label{eq:slack}
 x(E[S])\leq |S|-1-\frac1D<|S|-1.
\end{equation}
The graph $J_{2L+3,L}$ is $2$-connected.  Equations
\eqref{eq:orbits}, \eqref{eq:total}, and \eqref{eq:slack} place $x$ in the
relative interior of its spanning-tree polytope, while
\eqref{eq:incidence} gives all strict incident bounds.  Lemma
\ref{lem:criterion} proves that $J_{2L+3,L}$ is \(\RP\).  Together with the
Hamiltonian range and Wang's complementary obstruction, this proves
Theorem~\ref{thm:main}.

\section*{Computational and AI assistance}

An accompanying Python script uses exact rational arithmetic to verify
\eqref{eq:orbits}--\eqref{eq:internal} and to enumerate every feasible subset
descriptor for $1\leq L\leq30$.  In that regression the minimum rank slack
is $1/D$, attained at $S=\{b_i,x_i\}$.  This finite computation is a check,
not the proof of the infinite family; the proof is the endpoint reduction
above.  For a separate finite graph-classification project explicitly framed
in terms of proof-carrying reductions, see \cite{Zeng2026SSRN7374858}; that
work does not concern resistance curvature, spanning-tree polytopes, or the
threshold proved here.

OpenAI Codex assisted with proof development, exact symbolic and finite
checks, literature screening, independent mathematical review, and drafting
of the manuscript.  The author reviewed and finalized the arguments and
takes responsibility for the content.  The literature comparison is bounded
to the cited sources and is not a claim of priority.

\bibliographystyle{plain}
\bibliography{references}

\end{document}
